\section*{Abstract}
		Die T0-Time-Mass-Dualitäts-Theorie leitet fundamentale Konstanten und Massen parameterfrei aus dem universellen geometrischen Parameter $\xi = 4/30000$ ab. Dieses komplementäre Dokument prüft eine Validierung der Fraktaldimension $D_f = 3 - \xi \approx 2.99987$ durch Rückwärtsableitung aus dem experimentellen Massenverhältnis $r = m_{\mu} / m_e \approx 206.768$ (CODATA 2025). Die hier verwendete direkte geometrische Rechenkette trägt in der vorliegenden Form nicht: Die Elektronformel ergibt $0{,}0854$~GeV statt $m_e$, die Myon-Kette ergibt $m_e$ statt $m_\mu$, und die Nambu-Umkehrung ist mit der Torsionsmasse $m_T$ nicht konsistent; eine Validierung von $D_f$ folgt daraus nicht. Die systematische Massenberechnung präsentiert \emph{006\_T0\_Teilchenmassen\_De.pdf}.

	
	
	\section{Einleitung}
	\label{sec:einfuehrung}
	
	\begin{important}{Dokumenten-Komplementarität}{}
		Dieses Dokument konzentriert sich auf die \textbf{Prüfung einer Validierung der Fraktaldimension} $D_f$ aus experimentellen Lepton-Massen. Es ergänzt das Hauptdokument \emph{006\_T0\_Teilchenmassen\_De.pdf}, das die vollständige systematische Massenberechnung für alle Fermionen präsentiert.
	\end{important}
	
	Die Teilchenphysik steht vor dem fundamentalen Problem willkürlicher Massenparameter im Standardmodell. Die T0-Time-Mass-Dualitäts-Theorie revolutioniert diesen Ansatz durch eine vollständig parameterfreie Beschreibung.
	
	\section{Parameter und Grundformeln}
	\label{sec:parameter}
	
	Die Theorie basiert auf der Zeit-Energie-Dualität und fraktaler Raumzeit-Struktur.
	
	\subsection{Exakte geometrische Parameter}
	\label{subsec:exakte_parameter}
	
	\begin{align}
		\xi &= \frac{4}{30000} = \frac{1}{7500} \approx 1.333 \times 10^{-4}, \label{eq:xi} \\
		D_f &= 3 - \xi \approx 2.99986667, \label{eq:Df} \\
		\alpha &= \frac{1 - \xi}{137} \approx 7.298 \times 10^{-3}, \label{eq:alpha} \\
		K_{\text{frak}} &= 1 - 100 \xi \approx 0.9867, \label{eq:K} \\
		g_{T0}^2 &= \alpha K_{\text{frak}}, \label{eq:gT0} \\
		E_0 &= \frac{1}{\xi} \approx \SI{7500}{\giga\electronvolt}, \label{eq:E0} \\
		p &= -\frac{2}{3}. \label{eq:p}
	\end{align}

	$1/\xi = 7500$ ist eine reine Zahl; die Zuordnung der Einheit GeV in Gl.~\eqref{eq:E0} ist gesetzt, nicht hergeleitet (andere Dokumente lesen dieselbe Zahl als eV). In T0-Einheiten mit $\alpha = 1$ gilt $\xi E_0^2 = 1$, also $E_0^2 = 1/\xi = 7500$; die Zahl gehört zu $E_0^2$, nicht zu $E_0$ (Dok.~386). Der Korpuswert ist $E_0 = \sqrt{m_e m_\mu/K_{\text{frak}}} = 7{,}398$~MeV; $7500$~GeV liegt um den Faktor $1{,}01\times10^{6}$ darüber. Die folgenden Formeln verwenden $E_0\,\xi = 1$~GeV und hängen von dieser Setzung ab.
	
	\begin{result}{Präzision der Feinstrukturkonstante}{}
		Die Abweichung von $\alpha$ zu CODATA beträgt nur $\approx 0.013\%$ -- ein starkes Indiz für die fraktale Korrektur.
	\end{result}
	
	\section{Geometrische Ableitung der Massen - Direkte Methode}
	\label{sec:geometrische_ableitung}
	
	Die T0-Theorie bietet mehrere mathematisch äquivalente Methoden zur Massenberechnung. In diesem Dokument wird die \textbf{direkte geometrische Methode} daraufhin geprüft, ob sie die Fraktaldimension validiert.
	
	\subsection{Elektron-Masse $m_e$ - Direkte geometrische Methode}
	\label{subsec:elektron_masse}
	
	In der direkten geometrischen Methode:
	\begin{align}
		m_e &= E_0 \cdot \xi \cdot \sqrt{\alpha} \cdot \frac{\Gamma(D_f)}{\Gamma(3)}. \label{eq:me_direct}
	\end{align}

	Mit den Parametern aus Abschnitt~\ref{subsec:exakte_parameter} ($E_0\,\xi = 1$~GeV, $\sqrt{\alpha} = 0{,}08543$, $\Gamma(2{,}99987)/\Gamma(3) = 0{,}99994$) ergibt die Formel $m_e = 0{,}0854$~GeV -- um den Faktor $167$ größer als der Zielwert $5{,}10\times10^{-4}$~GeV. Die Formel liefert $m_e$ in dieser Form nicht. Die Abweichung $-0{,}20\,\%$ zu CODATA ($\SI{0.000511}{\giga\electronvolt}$) und die folgende Tabelle beziehen sich auf den Zielwert, nicht auf das Ergebnis der Formel.
	
	\subsection{Konsistenz-Check mit Hauptdokument}
	\label{subsec:konsistenz_check}
	
	\begin{table}[H]
		\centering
		\adjustbox{max width=\textwidth}{%
\begin{tabular}{lccc}
			\toprule
			\textbf{Methode} & \textbf{$m_e$ [GeV]} & \textbf{Genauigkeit} & \textbf{Quelle} \\
			\midrule
			Direkte geometrische (Zielwert) & $5.10\times10^{-4}$ & $99.8\%$ & Dieses Dokument \\
			Erweiterte Yukawa & $5.11\times10^{-4}$ & $99.9\%$ & 006\\
			Experiment (CODATA) & $5.11\times10^{-4}$ & $100\%$ & Referenz \\
			\bottomrule
		\end{tabular}%
}
		\caption{Konsistenz der Massenberechnungsmethoden in der T0-Theorie}
		\label{tab:methoden_konsistenz}
	\end{table}
	
	\begin{result}{Methodenvergleich}{}
		Die Übereinstimmung innerhalb von $0.2\%$ gilt für den Zielwert der direkten Methode, nicht für das Ergebnis von Gl.~\eqref{eq:me_direct}; eine Äquivalenz der Methoden ist damit nicht gezeigt. Die Yukawa-Methode schlägt die Brücke zum Standardmodell.
	\end{result}
	
	\subsection{Effektive Torsions-Masse $m_T$}
	\label{subsec:torsions_masse}
	
	\begin{align}
		R_f &= \frac{\Gamma(D_f)}{\Gamma(3)} \sqrt{\frac{E_0}{m_e}}, \label{eq:Rf} \\
		m_T &= \frac{m_e}{\xi} \sin(\pi \xi) \, \pi^2 \sqrt{\frac{\alpha}{K_{\text{frak}}}} \, R_f \approx \SI{5.220}{\giga\electronvolt}. \label{eq:mT}
	\end{align}
	
	\subsection{Myon-Masse $m_{\mu}$}
	\label{subsec:myon_masse}
	
	Aus RG-Dualität und Schleifenintegral $I$:
	\begin{align}
		I &= \int_0^1 \frac{m_e^2 x (1-x)^2}{m_e^2 x^2 + m_T^2 (1-x)}  dx \approx 1.60 \times 10^{-9}, \label{eq:I} \\
		r &\approx \sqrt{6 I} \approx 9.8 \times 10^{-5}, \label{eq:r} \\
		m_T \cdot r &\approx \SI{5.11e-4}{\giga\electronvolt}. \label{eq:mmu}
	\end{align}

	Das Integral ist numerisch mit $m_e = 5{,}11\times10^{-4}$~GeV und $m_T = 5{,}220$~GeV ausgewertet. Das Produkt $m_T\,r$ ergibt damit $m_e$, nicht $m_\mu$ ($\SI{0.105658}{\giga\electronvolt}$); die Kette liefert $m_\mu$ in dieser Form nicht.
	
	\begin{important}{Massenverhältnis}{}
		Da die Kette $m_\mu$ nicht liefert, lässt sich aus ihr kein Massenverhältnis $r = m_{\mu} / m_e$ berechnen; eine unabhängige Validierung der geometrischen Fundierung ergibt sich hier nicht.
	\end{important}
	
	\section{Rückwärts-Validierung: $D_f$ aus $r$ und Nambu-Formel}
	\label{sec:rueckwaerts_validierung}
	
	Die klassische Nambu-Formel $r \approx (3/2)/\alpha$ (Abw. $-0.58\%$) wird durch die $\xi$-Korrektur präzisiert.
	
	\subsection{Nambu-Umkehrung}
	\label{subsec:nambu_umkehrung}
	
	\begin{align}
		m_T^{\text{target}} &= \frac{m_{\mu}}{\sqrt{\alpha} \cdot (3/2) \cdot (1 - \xi)} = \frac{0{,}10566}{0{,}08543\times1{,}5\times0{,}99987} \approx \SI{0.825}{\giga\electronvolt}. \label{eq:mTtarget}
	\end{align}

	Dieser Wert ist mit Gleichung~\eqref{eq:mT} ($m_T = 5{,}220$~GeV) nicht konsistent; die Umkehrung trägt in dieser Form nicht.
	
	\subsection{Optimierung für $D_f$}
	\label{subsec:optimierung_df}
	
	Definiere $m_T(D_f)$ gemäß Gleichung~\ref{eq:mT} und löse:
	\begin{align}
		D_f = \arg\min \left| m_T(D_f) - m_T^{\text{target}} \right|. \label{eq:optDf}
	\end{align}
	
	\begin{keyresult}{Ergebnis der Rückwärts-Validierung}{}
		Da $m_T^{\text{target}} = 0{,}825$~GeV und $m_T = 5{,}220$~GeV nicht übereinstimmen, trägt die Optimierung auf dieser Grundlage nicht; ein Ergebnis $D_f \approx 2.99986667$ ist durch diese Rechnung nicht gestützt. Dass das experimentelle Massenverhältnis die fraktale Geometrie erzwingt, folgt aus dieser Rechnung nicht; die Grundlagen von \emph{006\_T0\_Teilchenmassen\_De.pdf} werden hier nicht unabhängig validiert.
	\end{keyresult}
	
	\section{Anwendung: Anomaler magnetischer Moment $a_{\mu}^{\text{T0}}$}
	\label{sec:anwendung_g2}
	
	Mit der Fraktaldimension $D_f = 3 - \xi$ und geometrischen Massen:
	\begin{align}
		F_2^{\text{T0}}(0) &= \frac{g_{T0}^2}{8 \pi^2} I_{\mu} K_{\text{frak}}, \label{eq:F2} \\
		\text{term} &= \left( \frac{\xi E_0}{m_T} \right)^p = m_T^{2/3}, \label{eq:term} \\
		F_{\text{dual}} &= \frac{1}{1 + \text{term}} \approx 0.249, \label{eq:Fdual} \\
		a_{\mu}^{\text{T0}} &= F_2^{\text{T0}}(0) \cdot F_{\text{dual}} \approx 1.53 \times 10^{-9} = 153 \times 10^{-11}. \label{eq:amu}
	\end{align}
	
	\begin{result}{Experimentelle Validierung}{}
		Stand 2025 (Fermilab-Endergebnis, White Paper 2025 der Muon~$g-2$ Theory Initiative): $\Delta a_\mu = (38 \pm 63) \times 10^{-11}$. Der T0-Wert $153 \times 10^{-11}$ liegt $(153 - 38)/63 \approx 1.8\sigma$ darüber.
	\end{result}
	
	\section{Python-Implementierung und Reproduzierbarkeit}
	\label{sec:python_implementierung}
	
	\begin{important}{Volle Transparenz}{}
		Zur Reproduktion aller numerischen Berechnungen siehe das externe Skript \texttt{t0\_df\_from\_masses\_geometry.py} im Repository-Ordner.
	\end{important}
	
	\section{Referenzen}
	\label{sec:referenzen}
	
	\begin{itemize}
		\item Pascher, J. (2025). \emph{T0-Modell: Vollständige parameterfreie Teilchenmassen-Berechnung} (006\_T0\_Teilchenmassen\_De.pdf). Verfügbar unter:\\ 
		
		\item Pascher, J. (2025). \emph{T0-Time-Mass-Duality Repository}, GitHub v1.6. Verfügbar unter: 
		
		\item CODATA (2025). \emph{Fundamentale physikalische Konstanten}, NIST.
	\end{itemize}
	